\documentclass[11pt]{article}

\usepackage[margin=1in]{geometry}
\usepackage{amsmath,amssymb,amsthm}
\usepackage{booktabs}
\usepackage{natbib}

\newtheorem{theorem}{Theorem}
\newtheorem{proposition}{Proposition}
\newtheorem{remark}{Remark}
\newtheorem{definition}{Definition}

\title{Predictive e-diagnostics for multi-state movement models\\A theoretical setup for hidden Markov models}
\author{Author information to be completed}
\date{Draft version: 2026-05-07}

\begin{document}
\maketitle

\begin{abstract}
Hidden Markov models are widely used to describe animal movement as a sequence of unobserved behavioural states that generate observed step lengths and turning angles. Standard model assessment workflows often combine likelihood-based criteria, decoded states, simulation checks and graphical diagnostics. These tools are useful, but they do not by themselves provide a sequentially valid account of adaptive model checking. This paper develops the theoretical core of a predictive e-diagnostic framework for multi-state movement models. A fitted HMM is treated as a sequential generator, and diagnostic alternatives are encoded through predictive density ratios computed on a validation sequence. The latent state process is integrated out by filtering, so the resulting diagnostics are based on observable predictive distributions rather than decoded state paths. Under a fixed train/test protocol, the cumulative product of predictive density ratios is a nonnegative martingale under the null HMM and hence an e-process. This yields anytime-valid monitoring of predictive evidence against the fitted model and provides a foundation for movement-specific diagnostics targeting state number, angular misspecification, step-angle dependence and state-duration structure.
\end{abstract}

\paragraph{Keywords.}
Animal movement; hidden Markov model; e-value; e-process; sequential prediction; model diagnostics; filtering.

\section{Introduction}

Hidden Markov models (HMMs) are a standard tool for analysing animal movement data when movement is viewed as the observable output of an unobserved behavioural state process. In typical applications, observations are constructed from locations as step lengths and turning angles, and the latent states are interpreted as behavioural modes such as resting, foraging or travelling. Software such as \texttt{moveHMM} \citep{michelot_2016_movehmm} and \texttt{momentuHMM} \citep{mcclintock_2018_momentuhmm} has made these models accessible for ecological analyses.

Model assessment for movement HMMs often relies on AIC or BIC, likelihood comparisons, inspection of decoded states, residual or simulation-based checks, and ecological interpretation of inferred behavioural states. These practices are valuable, but they raise two statistical issues. First, diagnostics are often chosen after seeing the data, which makes the inferential status of repeated exploratory checking unclear. Second, the latent states are sometimes treated as if they were observed after decoding, even though they remain model-dependent reconstructions.

This note proposes a predictive alternative. The fitted HMM is evaluated as a sequential generator of the validation trajectory. At each time $t$, the null model provides an observable predictive density $p_0(y_t \mid \mathcal F_{t-1})$. A diagnostic alternative provides another predictive density $p_t^1(y_t \mid \mathcal F_{t-1})$, chosen before $Y_t$ is observed. The ratio of these densities is a predictive e-value increment. Its cumulative product is an e-process under the null model, connecting movement model diagnostics to the modern literature on e-values and anytime-valid inference \citep{vovk_wang_2021_evalues,howard_2020_time_uniform}.

The contribution of this paper is deliberately narrow. It establishes the train/test version of the theory for fixed HMMs and fixed or predictable diagnostic alternatives. Extensions to cross-fitting, fully adaptive diagnostic selection and same-sample exploratory workflows are left for later stages of the project.

\section{Movement HMM as a sequential generator}

Consider a validation trajectory indexed by $t = 1, \ldots, T$. The observed movement variable is
\begin{equation}
Y_t = (L_t, \Theta_t),
\end{equation}
where $L_t$ is the step length and $\Theta_t$ is the turning angle. The latent state is
\begin{equation}
S_t \in \{1, \ldots, K\}.
\end{equation}
The observable filtration is
\begin{equation}
\mathcal F_t = \sigma(Y_1, \ldots, Y_t),
\end{equation}
with $\mathcal F_0$ containing only information fixed before validation, including any fitted model parameters from the training sample.

\begin{table}[h]
\centering
\begin{tabular}{ll}
\toprule
Notation & Meaning \\
\midrule
$Y_t$ & Observed movement variable at time $t$ \\
$L_t$ & Step length at time $t$ \\
$\Theta_t$ & Turning angle at time $t$ \\
$S_t$ & Latent behavioural state at time $t$ \\
$K$ & Number of states in the null HMM \\
$M_0$ & Null HMM, fixed before validation \\
$M_1$ & Diagnostic alternative, fixed or predictable before $Y_t$ \\
$p_0(y_t \mid \mathcal F_{t-1})$ & Observable predictive density under $M_0$ \\
$p_t^1(y_t \mid \mathcal F_{t-1})$ & Diagnostic predictive density at time $t$ \\
$E_t$ & Predictive e-value increment at time $t$ \\
$E_{1:t}$ & Cumulative product of increments up to time $t$ \\
$\alpha$ & Anytime-valid monitoring level \\
\bottomrule
\end{tabular}
\caption{Notation used in the theoretical setup.}
\label{tab:notation}
\end{table}

The null model $M_0$ is a $K$-state HMM with initial distribution $\boldsymbol\pi_0$, transition matrix $\boldsymbol\Gamma_0$ and state-dependent observation densities. Write
\begin{equation}
f_{0,s}(y_t) = p_0(y_t \mid S_t=s)
\end{equation}
for the observation density in state $s$. In a common movement specification,
\begin{equation}
f_{0,s}(l_t,\theta_t) = f_{0,s}^L(l_t) f_{0,s}^{\Theta}(\theta_t),
\end{equation}
for example with Gamma step lengths and von Mises turning angles. This conditional independence factorisation is a modelling choice, not a requirement of the e-process construction.

\section{Observable predictive density}

The key object is the one-step-ahead observable predictive density. Since $S_t$ is latent, this density must marginalize over the predictive distribution of the state:
\begin{equation}
\label{eq:p0_predictive}
p_0(y_t \mid \mathcal F_{t-1})
=
\sum_{s=1}^K
f_{0,s}(y_t)\Pr_0(S_t=s \mid \mathcal F_{t-1}).
\end{equation}
This is the density of the next observed movement variable under the fitted HMM, conditional only on the observed past.

Let
\begin{equation}
a_{t\mid t-1}(s) = \Pr_0(S_t=s \mid \mathcal F_{t-1})
\end{equation}
denote the predictive state probabilities before observing $Y_t$. Then
\begin{equation}
p_0(y_t \mid \mathcal F_{t-1}) = \sum_{s=1}^K a_{t\mid t-1}(s) f_{0,s}(y_t).
\end{equation}
After observing $Y_t=y_t$, the filtered probabilities are
\begin{equation}
\label{eq:filter_update}
a_{t\mid t}(s)
=
\frac{a_{t\mid t-1}(s) f_{0,s}(y_t)}
{\sum_{r=1}^K a_{t\mid t-1}(r) f_{0,r}(y_t)}.
\end{equation}
The next predictive state probabilities are
\begin{equation}
\label{eq:filter_predict}
a_{t+1\mid t}(s) = \sum_{r=1}^K a_{t\mid t}(r)\Gamma_0(r,s).
\end{equation}
Equations \eqref{eq:p0_predictive} to \eqref{eq:filter_predict} are the computational basis of the null predictive density.

\section{Diagnostic alternatives and predictive e-values}

A diagnostic alternative is represented by a predictive density
\begin{equation}
p_t^1(y_t \mid \mathcal F_{t-1}).
\end{equation}
The subscript $t$ allows the alternative to depend on the observed past. The important requirement is predictability: the rule defining $p_t^1$ must be fixed before observing $Y_t$. This includes alternatives estimated on a separate training sample and alternatives updated using only $\mathcal F_{t-1}$.

Examples include:
\begin{itemize}
\item a $K+1$ state HMM used to diagnose whether the null HMM has too few states;
\item a more flexible angular distribution used to diagnose misspecified turning angles;
\item a joint step-angle model used to diagnose residual dependence between $L_t$ and $\Theta_t$ within states;
\item a duration-aware alternative used to diagnose non-geometric state dwell times.
\end{itemize}
These alternatives are diagnostic bets, not necessarily claims about the true ecological data-generating process.

\begin{definition}[Predictive e-value increment]
For $t=1,\ldots,T$, define
\begin{equation}
\label{eq:increment}
E_t =
\frac{p_t^1(Y_t \mid \mathcal F_{t-1})}
{p_0(Y_t \mid \mathcal F_{t-1})},
\end{equation}
whenever the ratio is well-defined under $M_0$.
\end{definition}

\begin{definition}[Predictive e-process]
The cumulative process is
\begin{equation}
\label{eq:eprocess}
E_{1:t} = \prod_{u=1}^t E_u,
\qquad E_{1:0}=1.
\end{equation}
Equivalently,
\begin{equation}
\log E_{1:t}
=
\sum_{u=1}^t
\left\{
\log p_u^1(Y_u \mid \mathcal F_{u-1})
-
\log p_0(Y_u \mid \mathcal F_{u-1})
\right\}.
\end{equation}
\end{definition}

\section{Main theoretical guarantee}

\begin{theorem}[Predictive e-process for a fixed HMM]
\label{thm:predictive_eprocess}
Assume that the null model $M_0$ defines the observable predictive density $p_0(y_t \mid \mathcal F_{t-1})$ at each time $t$. Assume that the diagnostic predictive density $p_t^1(y_t \mid \mathcal F_{t-1})$ is chosen before observing $Y_t$ and is a valid density. If the validation sequence is generated under $M_0$, then
\begin{equation}
\mathbb E_0(E_t \mid \mathcal F_{t-1}) \leq 1.
\end{equation}
If $p_t^1(\cdot \mid \mathcal F_{t-1})$ is absolutely continuous with respect to $p_0(\cdot \mid \mathcal F_{t-1})$ and integrates to one on the support of $p_0$, then equality holds. Consequently, $(E_{1:t})_{t\geq0}$ is a nonnegative supermartingale under $M_0$, and it is a martingale under the equality condition.
\end{theorem}

\begin{proof}
Conditional on $\mathcal F_{t-1}$, the validation observation $Y_t$ has density $p_0(y \mid \mathcal F_{t-1})$ under $M_0$. Therefore,
\begin{align}
\mathbb E_0(E_t \mid \mathcal F_{t-1})
&=
\int
\frac{p_t^1(y \mid \mathcal F_{t-1})}
{p_0(y \mid \mathcal F_{t-1})}
 p_0(y \mid \mathcal F_{t-1})\,dy \\
&=
\int_{\{p_0>0\}} p_t^1(y \mid \mathcal F_{t-1})\,dy \\
&\leq 1.
\end{align}
The inequality becomes equality when the diagnostic density assigns all its mass to the support of the null predictive density. For the cumulative process,
\begin{align}
\mathbb E_0(E_{1:t} \mid \mathcal F_{t-1})
&=
E_{1:t-1}\mathbb E_0(E_t \mid \mathcal F_{t-1}) \\
&\leq E_{1:t-1}.
\end{align}
Thus $(E_{1:t})_{t\geq0}$ is a nonnegative supermartingale. Under equality it is a martingale.
\end{proof}

\begin{proposition}[Anytime-valid monitoring]
\label{prop:anytime}
Under the conditions of Theorem \ref{thm:predictive_eprocess}, for any $\alpha \in (0,1)$,
\begin{equation}
\mathbb P_0\left(\sup_{t\geq0} E_{1:t} \geq \frac{1}{\alpha}\right) \leq \alpha.
\end{equation}
Equivalently,
\begin{equation}
\mathbb P_0\left(\sup_{t\geq0} \log E_{1:t} \geq \log(1/\alpha)\right) \leq \alpha.
\end{equation}
\end{proposition}

\begin{proof}
This is Ville's maximal inequality applied to the nonnegative supermartingale $(E_{1:t})_{t\geq0}$ with initial value one. This is the standard mechanism underlying anytime-valid inference with e-processes \citep{howard_2020_time_uniform}.
\end{proof}

\section{Why filtering is essential}

\begin{proposition}[Validity is based on marginal observable prediction]
\label{prop:latent_filtering}
For an HMM, the denominator in \eqref{eq:increment} must be the marginal observable predictive density
\begin{equation}
p_0(y_t \mid \mathcal F_{t-1})
=
\sum_{s=1}^K f_{0,s}(y_t)\Pr_0(S_t=s \mid \mathcal F_{t-1}).
\end{equation}
The e-process validity does not require observing $S_t$ and does not follow from conditioning on a decoded state path $\widehat S_t$.
\end{proposition}

\begin{proof}
The proof of Theorem \ref{thm:predictive_eprocess} only requires that the denominator be the correct conditional density of the observable variable $Y_t$ given $\mathcal F_{t-1}$. In an HMM, this density is obtained by summing over latent states using the predictive state probabilities. Decoded state paths do not appear in this conditional law.
\end{proof}

\begin{remark}[Decoded states]
A Viterbi path or any other decoded state sequence is a model-dependent reconstruction, not an observed state sequence. Decoded states may be useful for descriptive interpretation after a diagnostic signal has been found. They should not be used as if $S_t$ were observed when establishing the e-process property.
\end{remark}

\section{Estimated parameters and validation design}

The cleanest version of the guarantee is conditional on a training set. Let $\mathcal D_{\mathrm{train}}$ denote the data used to fit the null HMM and, if applicable, the diagnostic alternatives. The validation sequence is then used only to compute predictive density ratios. Conditional on $\mathcal D_{\mathrm{train}}$, the fitted parameters are fixed, and Theorem \ref{thm:predictive_eprocess} applies to the resulting predictive densities.

This train/test design is the recommended starting point for the first article because it gives reviewers a transparent route from fitted movement models to anytime-valid diagnostic monitoring. A multi-individual design is particularly attractive in movement ecology: the null model and diagnostic alternatives can be fitted on some individuals and evaluated on held-out individuals.

If the same observations are used to select the model, estimate its parameters, choose the diagnostics and compute the e-process, the exact guarantee no longer follows directly from Theorem \ref{thm:predictive_eprocess}. Such workflows may still be useful as exploratory diagnostics, but they should be labelled as exploratory. Later extensions may use cross-fitting, pre-specified diagnostic menus or conditionally valid adaptive choices.

\section{Diagnostic interpretation}

Large values of $E_{1:t}$ indicate cumulative predictive evidence against the null HMM relative to the diagnostic alternative used to define $p_t^1$. The interpretation is comparative and diagnostic. It does not imply that the alternative is the true data-generating mechanism.

Under a data-generating process $P^\star$, the log e-process is
\begin{equation}
\log E_{1:t}
=
\sum_{u=1}^t
\log
\frac{p_u^1(Y_u \mid \mathcal F_{u-1})}
{p_0(Y_u \mid \mathcal F_{u-1})}.
\end{equation}
If the diagnostic alternative predicts better on average, in the sense that the expected log ratio is positive over relevant segments of the trajectory, then $\log E_{1:t}$ tends to increase over those segments. If the alternative targets the wrong defect, the e-process may remain insensitive even when the null HMM is ecologically inadequate.

For movement applications, this suggests reporting not only whether the e-process crosses the threshold $\log(1/\alpha)$, but also where the local increments
\begin{equation}
\log E_t = \log p_t^1(Y_t \mid \mathcal F_{t-1}) - \log p_0(Y_t \mid \mathcal F_{t-1})
\end{equation}
are large. This localizes model failure in time, by individual, by habitat or by biologically meaningful period.

\section{Discussion and next steps}

This theoretical setup supports a first article with a focused scope. The null HMM is assessed as a sequential generator of validation data. The latent states are integrated out through filtering. Diagnostic alternatives define targeted predictive bets. The resulting cumulative ratios are e-processes under the fixed train/test null, giving anytime-valid thresholds for sequential model criticism.

The next computational step is to implement a minimal R prototype that simulates a simple HMM, computes observable predictive densities by filtering and verifies the null behaviour of the e-process. Subsequent diagnostics can then be developed for missing states, angular misspecification, residual step-angle dependence and duration misspecification.

\bibliographystyle{plainnat}
\bibliography{references}

\end{document}
